\section{Radial equality and the nonlinear conclusion}
\label{ads:sec:conclusion}

\begin{proposition}[Exact radial equality]
\label{ads:prop:radial}
Suppose the fixed TT seed $q$ is invariant under rotations of $S^2$.
Every branch in Theorem~\ref{ads:thm:main} is radial for all sufficiently
small $\eps$, and on that parameter interval
\[
 m_{\AH}(\eps)=b_*(A_\eps).
\]
This includes the zero seed, and requires no sign for $q_{ss}|_S$.
\end{proposition}

\begin{proof}
We first prove radiality of the actual solution, rather than merely
of its Taylor coefficients.  Fix a rotation $\mathcal R$, and set
$\widehat\phi=\phi_\eps\circ\mathcal R$ and
$w=\widehat\phi-\phi_\eps$.
Because $q$ is radial, $\widehat\phi$ satisfies the same equation and
boundary condition as $\phi_\eps$.  Their difference satisfies
\begin{equation}
 (\Delta_{g_0}-c_\eps)w=0,
 \qquad w_s|_S=c_\eps^*w|_S,
 \label{ads:eq:radial-uniqueness-equation}
\end{equation}
where divided differences of the nonlinearities give, uniformly on
$M$ and uniformly in $\mathcal R$,
\[
 c_\eps=3+o(1),\qquad c_\eps^*=O(\eps).
\]
Indeed, the derivatives of the interior and boundary nonlinearities
with respect to their positive scalar argument $z$ are respectively
\[
 \frac34(5z^4-1)+\frac{7\eps^2}{8}|q|^2z^{-8},
 \qquad -\frac{3\eps}{4}q_{ss}z^{-4}.
\]
The asserted bounds follow from boundedness of $q$ and the given
uniform convergence $\phi_\eps\to1$.

Choose $3/2<\beta<\min(\tau,\sqrt3)$. Then
$w=o(e^{-\beta s})$, and \eqref{ads:eq:radial-uniqueness-equation} has the
form of Lemma~\ref{ads:lem:inverse} with $V=c_\eps-3$, Robin coefficient
$a_*=c_\eps^*$, and zero interior and boundary data. For sufficiently
small $\eps$, uniformly in the rotation, $\|V\|_\infty\leq c_\beta$
and $\|a_*\|_\infty\leq\beta/2$. The comparison estimate gives $w=0$,
so the actual $\phi_\eps$ is radial.

Write the invariant seed as
$q=A_q(s)\,\dd s^2+B_q(s)r^2\sigma$.
Its trace and divergence equations give
\[
 B_q=-A_q/2,
 \qquad A_q'+3(N/r)A_q=0.
\]
Consequently, for a constant $C_q\in\mathbb R$,
\begin{equation}
 q=C_qr^{-3}\left(\dd s^2-\frac12r^2\sigma\right).
 \label{ads:eq:radial-tt-tensor}
\end{equation}
Fix a sufficiently small parameter.  Define proper radial distance
$\ell$ and the actual area radius $R$ by
\[
 \dd\ell=\phi_\eps^2\,\dd s,
 \qquad R=r\phi_\eps^2.
\]
Then the actual metric and second-form data are
\[
 g_\eps=\dd\ell^2+R^2\sigma,
 \qquad K_\eps=-2k\,\dd\ell^2+kR^2\sigma,
 \qquad k=-\frac{\eps C_q}{2R^3}.
\]
A dot denotes an $\ell$-derivative.  The momentum and scalar constraints
are therefore
\begin{equation}
 \dot k=-3\frac{\dot R}{R}k,
 \qquad
 -\frac{4\ddot R}{R}
 +\frac{2(1-\dot R^2)}{R^2}+6=6k^2.
 \label{ads:eq:radial-constraints}
\end{equation}
The scalar-curvature formula here follows from the sectional curvatures
$-\ddot R/R$ and $(1-\dot R^2)/R^2$ of this warped product.

Define
\begin{equation}
 \mathcal M
 =\frac R2\left(1+R^2-\dot R^2+R^2k^2\right).
 \label{ads:eq:radial-conserved-mass}
\end{equation}
Direct differentiation, using the momentum equation, gives
\[
 \dot{\mathcal M}
 =\frac{\dot R}{2}
 \left(1+3R^2-\dot R^2-2R\ddot R-3R^2k^2\right)=0
\]
by the scalar equation in \eqref{ads:eq:radial-constraints}.
No division by $\dot R$ has been used, so conservation is valid at
turning points of $R$ as well.
At the boundary the MOTS equation reads
\[
 \frac{2\dot R}{R}+2k=0.
\]
Substitution in \eqref{ads:eq:radial-conserved-mass} yields
\begin{equation}
 \mathcal M|_S
 =\frac{R|_S}{2}\bigl(1+(R|_S)^2\bigr)
 =b_*\bigl(4\pi(R|_S)^2\bigr)=b_*(A_\eps).
 \label{ads:eq:radial-mass-at-boundary}
\end{equation}

We identify the same constant with the specified mass flux at infinity.
Put $\psi=\phi_\eps-1$.  Its actual radial equation gives
\[
 \psi''+2\frac Nr\psi'-3\psi
 =O(\psi^2)+O(r^{-6}),
\]
where \eqref{ads:eq:radial-tt-tensor} was used in the source estimate.
The assumed weighted decay gives
$\psi,\psi'=O(e^{-\tau s})$, and $N/r=1+O(e^{-2s})$.
It follows that
\[
 \psi''+2\psi'-3\psi=O(e^{-\gamma s}),
 \qquad \gamma=\min(2\tau,\tau+2,6)>3.
\]
Variation of constants for the roots $1,-3$ excludes the growing
solution by the prescribed decay and yields, as in
Lemma~\ref{ads:lem:inverse}, a constant $c=c(\eps)$ with
\begin{equation}
 \phi_\eps=1+cr^{-3}+o(r^{-3}),
 \qquad
 \partial_s\phi_\eps=-3cNr^{-4}+o(r^{-3}).
 \label{ads:eq:radial-actual-asymptotic}
\end{equation}
In particular,
\[
 R=r+2cr^{-2}+o(r^{-2}),\qquad
 \dot R=N+2r\frac{\partial_s\phi_\eps}{\phi_\eps}
 =N-6cNr^{-3}+o(r^{-2}).
\]
Using $N^2=1+r^2-2m_0/r$, we obtain
\[
 1+R^2-\dot R^2
 =\frac{2m_0+16c}{r}+o(r^{-1}).
\]
Also $R^3k^2=O(R^{-3})$, so
\begin{equation}
 \lim_{\ell\to\infty}\mathcal M=m_0+8c.
 \label{ads:eq:radial-mass-at-infinity}
\end{equation}

For a direct flux normalization check, rotational invariance implies
$p_i=0$ for $i=1,2,3$.
Equation~\eqref{ads:eq:radial-actual-asymptotic} implies
\[
 g_\eps-g_0=4cr^{-3}b+o(r^{-3})
\]
in scaled components, with its first differentiated asymptotic.
For a perturbation $fb$ in dimension three, the mass-flux covector
integrand with test function $V$ simplifies to
$2(f\,\dd V-V\,\dd f)$.
With $f=4cr^{-3}$, $V=V_0=\sqrt{1+r^2}$, and
$\nu_b=\sqrt{1+r^2}\,\partial_r$, its normal component is
\[
 2\bigl(f\partial_{\nu_b}V_0
       -V_0\partial_{\nu_b}f\bigr)
 =32cr^{-2}+24cr^{-4}.
\]
The error terms give zero limiting flux.  Integrating over a coordinate
sphere of area $4\pi r^2$ and dividing by $16\pi$ gives
$p_0-m_0=8c$.
Since $p_0\to m_0>0$ as $\eps\to0$,
\[
 m_{\AH}=p_0=m_0+8c=\mathcal M=b_*(A_\eps)
\]
for sufficiently small $\eps$, by
\eqref{ads:eq:radial-mass-at-boundary}--\eqref{ads:eq:radial-mass-at-infinity}.
This argument includes $C_q=0$.  In that case one can also apply the
comparison estimate of Lemma~\ref{ads:lem:inverse} to $\phi_\eps-1$
and conclude $\phi_\eps\equiv1$ directly.
\end{proof}

\begin{proof}[Completion of the proof of Theorem~\ref{ads:thm:main}]
Fix the seed and the given solution branch.  If $q$ is radial,
Proposition~\ref{ads:prop:radial} gives exact equality on a sufficiently
small parameter interval.
Suppose henceforth that $q$ is not radial.
Proposition~\ref{ads:prop:quadratic} gives the actual expansion
\[
 \mathfrak D_q(\eps)=c_2[q]\eps^2+o(\eps^2),
 \qquad c_2[q]\geq0.
\]
If $c_2[q]>0$, this proves strict positivity of the deficit for every
sufficiently small positive $\eps$.

If $c_2[q]=0$, Corollary~\ref{ads:cor:kernel} gives
$q=Q(v)$ with $Lv=0$ and with only constant and linear angular modes.
Its linear part is nonzero, for otherwise $q$ would be radial.
These homogeneous modes have the stronger differentiated decay needed
for the actual fourth-order expansion in
Proposition~\ref{ads:prop:expansions}.
Proposition~\ref{ads:prop:quartic} then gives
\[
 \mathfrak D_q(\eps)=c_4[q]\eps^4+o(\eps^4),
 \qquad c_4[q]>0,
\]
with all coefficients below order four equal to zero.
Again the actual deficit is strictly positive for every sufficiently
small positive $\eps$.

At $\eps=0$ the background identity
$m_0=b_*(4\pi a^2)$ gives equality.
The preceding alternatives exhaust all fixed seeds.  Choose
$\eps_0>0$ small enough for the applicable alternative and the given
branch, with $\eps_0\leq\eps_*$.
This proves the exact inequality on $0\leq\eps<\eps_0$, equality for
radial seeds, and strict inequality for nonradial seeds at positive
parameters in that interval.
The size of the interval is allowed to depend on the fixed seed and
branch; no uniform positive lower bound for $c_2$ or $c_4$ is asserted
or required.
\end{proof}
